% careerdossier.tex — the CareerDossierTeX user manual.
%
% Copyright (C) 2026 Amir Sadeghi.
%
% This file is distributed with the Work under the LaTeX Project Public
% License, version 1.3c or, at your option, any later version. The Work is
% defined in manifest.txt, which also records this file under "Distributed
% with the Work".
%
% LPPL maintenance status: `maintained'.
% Current Maintainer: Amir Sadeghi.
%
% Build it with the repository target, which keeps every artifact out of the
% source tree:
%
%     make manual          -> build/manual/careerdossier.pdf
%
% The manual documents the released public interface only. Private LaTeX3
% names are deliberately absent: tests/lint/run-manual-names.sh fails the
% build if one appears here, and fails it again if a public name documented
% here is not declared by a file in the Work.
%
% The version and date below are the release this manual documents, and
% tests/lint/run-version-declarations.sh checks them against manifest.txt's
% Work files.

\documentclass[11pt]{article}

\newcommand{\CDossierManualVersion}{0.8.0}
\newcommand{\CDossierManualDate}{2026-08-12}

\usepackage[letterpaper,margin=1in]{geometry}
\usepackage{fontspec}
\usepackage{longtable}
\usepackage{booktabs}
\usepackage{array}
% Table cells hold long names; justifying them produces underfull lines that
% carry no information. L is p, set ragged right.
\newcolumntype{L}[1]{>{\raggedright\arraybackslash}p{#1}}
\usepackage{fancyvrb}
\usepackage{enumitem}
\usepackage[hidelinks]{hyperref}

\setmainfont{TeX Gyre Termes}
\setsansfont{TeX Gyre Heros}[Scale=MatchLowercase]
\setmonofont{TeX Gyre Cursor}[Scale=MatchLowercase]

\hypersetup{
  pdftitle          = {CareerDossierTeX: a LuaLaTeX toolkit for career documents},
  pdfauthor         = {Amir Sadeghi},
  pdfsubject        = {User manual for the CareerDossierTeX document classes},
  pdflang           = {en},
  pdfdisplaydoctitle = true,
  bookmarksnumbered = true
}

\DefineVerbatimEnvironment{latexcode}{Verbatim}
  {fontsize=\small, frame=leftline, framesep=8pt, xleftmargin=10pt}

\DefineVerbatimEnvironment{plaincode}{Verbatim}
  {fontsize=\small, xleftmargin=10pt}

% Public names are long, and to TeX each is one unbreakable token: a paragraph
% or a table cell holding \CDossierRecordListEdgeAboveSkip overruns the measure
% rather than wrapping. Break them where a reader already reads a boundary ---
% before an interior capital, and after `-', `/', or `.' --- and insert no
% hyphen, because a broken name must still read as the name it is. This is the
% same rule the toolkit applies to an address in \CDossierLink.
\makeatletter
\newif\ifcd@firstchar
\newcommand{\cd@breaklong}[1]{%
  \cd@firstchartrue
  \@tfor\cd@char:=#1\do{\expandafter\cd@testchar\cd@char}}
\newcommand{\cd@testchar}[1]{%
  \ifcd@firstchar
    \cd@firstcharfalse
  \else
    \ifnum\uccode`#1=`#1 \allowbreak\fi
  \fi
  #1%
  \ifnum0\if#1-1\fi\if#1/1\fi\if#1.1\fi>0 \allowbreak\fi}

% Robust, because these appear in section titles: the scanner above is built
% from \if tests, and an unprotected one expanded while writing the table of
% contents or a PDF bookmark ends the run with `Incomplete \iffalse'.
\DeclareRobustCommand{\opt}[1]{\texttt{\cd@breaklong{#1}}}
\DeclareRobustCommand{\key}[1]{\texttt{\cd@breaklong{#1}}}
\DeclareRobustCommand{\cmd}[1]{\texttt{\textbackslash\cd@breaklong{#1}}}
\DeclareRobustCommand{\env}[1]{\texttt{\cd@breaklong{#1}}}
\DeclareRobustCommand{\file}[1]{\texttt{\cd@breaklong{#1}}}
\makeatother

% A PDF bookmark is a string, not typeset material, so give hyperref a plain
% spelling of each rather than leaving it to warn about a command it cannot
% render.
\pdfstringdefDisableCommands{%
  \def\opt#1{#1}%
  \def\key#1{#1}%
  \def\cmd#1{\textbackslash #1}%
  \def\env#1{#1}%
  \def\file#1{#1}%
}

% An address in running text or a table cell gets the same treatment, through
% url.sty's own breakpoint list rather than the scanner above, since an address
% carries characters the scanner would have to be told about one at a time.
\makeatletter
\g@addto@macro\UrlBreaks{%
  \do\a\do\b\do\c\do\d\do\e\do\f\do\g\do\h\do\i\do\j\do\k\do\l\do\m%
  \do\n\do\o\do\p\do\q\do\r\do\s\do\t\do\u\do\v\do\w\do\x\do\y\do\z%
  \do\A\do\B\do\C\do\D\do\E\do\F\do\G\do\H\do\I\do\J\do\K\do\L\do\M%
  \do\N\do\O\do\P\do\Q\do\R\do\S\do\T\do\U\do\V\do\W\do\X\do\Y\do\Z%
  \do\0\do\1\do\2\do\3\do\4\do\5\do\6\do\7\do\8\do\9%
  \do\=\do\?\do\&\do\-\do\_}
\makeatother
\DeclareUrlCommand{\addr}{\urlstyle{tt}}

\setlist{noitemsep, topsep=4pt}
\setlength{\parskip}{6pt}
\setlength{\parindent}{0pt}

% A paragraph naming three public commands has few break points and would
% otherwise set overfull rather than merely loose.
\tolerance=800
\emergencystretch=1em

\title{\sffamily\bfseries CareerDossierTeX}
\author{Amir Sadeghi}
\date{Version \CDossierManualVersion\ (\CDossierManualDate)}

\begin{document}

\maketitle

\begin{center}
\begin{minipage}{0.86\textwidth}
\small
A reusable LuaLaTeX toolkit for producing a consistent set of career documents
--- résumé, cover letter, academic CV, and seven kinds of statement --- from one
shared profile.

Distributed under the \LaTeX{} Project Public License, version 1.3c or, at your
option, any later version. LPPL maintenance status \texttt{maintained}; current
maintainer Amir Sadeghi. The Work is defined in \file{manifest.txt}.

Repository and issue tracker:
\url{https://github.com/amirhs1/CareerDossierTeX}
\end{minipage}
\end{center}

\vspace{1.5\baselineskip}

\tableofcontents

\clearpage

\section{About this manual}

This manual documents the released public interface of CareerDossierTeX: every
class, option, key, command, environment, and design token an author writes,
with its accepted values and its default. It is the reference for people writing
documents with the toolkit.

It does not document the internals. Which module owns which concern, and why, is
recorded in \file{docs/ARCHITECTURE.md} in the repository. Private \LaTeX3 names
--- those spelled with a leading double underscore --- are not public and do not
appear here; they may be renamed or removed in any release.

\subsection{Interface stability}

Before \texttt{v0.10.0} the public interface may still change between minor
versions. Every such change is recorded in \file{CHANGELOG.md}, and a change
that needs an edit to your source is additionally recorded, with a before-and-after
example, in \file{docs/MIGRATION.md}. Read the latter before upgrading.

Calibrated token \emph{values} are not stable before \texttt{v0.10.0}; the token
names and the boundaries they own are.

\subsection{What is supported}

Only behaviour documented in this manual, and covered by the repository's tests
and shipped examples, should be treated as supported. Internal helper commands
are not public merely because they are technically reachable.

Text extraction is a tested baseline, not an accessibility guarantee. This
toolkit makes no PDF/UA, WCAG, or ATS conformance claim; \S\ref{sec:tagging}
states the scope of what has actually been verified.

\section{Requirements and installation}
\label{sec:install}

\subsection{Requirements}

\begin{itemize}
  \item \textbf{LuaLaTeX}, which is the sole supported engine. XeLaTeX and
        pdf\LaTeX{} stop with a fatal error naming LuaLaTeX; there is no
        compatibility mode and no way to bypass the guard.
  \item A current \TeX{} Live or MiK\TeX{} distribution providing the
        \textsf{TeX Gyre} fonts (Termes and Heros), \textsf{fontspec},
        \textsf{geometry}, \textsf{hyperref}, \textsf{enumitem}, and
        \textsf{xcolor}. Nothing is bundled with this toolkit.
  \item \textbf{latexmk}, recommended: a résumé or CV needs two passes, and
        latexmk performs them.
  \item \textbf{Biber and BibLaTeX}, only if you opt into the bibliography
        integration of \S\ref{sec:biblatex}.
\end{itemize}

\subsection{Three ways to install}

\subsubsection*{Route 1 --- beside the document}

Copy the \file{.sty} and \file{.cls} files into the same directory as your
document. \LaTeX{} finds files in the current directory first, so nothing else
is needed. This is the right route for a single application, for a directory you
will send to someone else, and for trying the toolkit out.

\subsubsection*{Route 2 --- a local \texttt{texmf} tree}

Install once and use the classes from any directory:

\begin{plaincode}
mkdir -p ~/texmf/tex/latex/careerdossier
cp careerdossier-*.sty careerdossier-*.cls ~/texmf/tex/latex/careerdossier/
\end{plaincode}

On \TeX{} Live the tree is found without a further step; \texttt{mktexlsr} is
needed only for a system-wide tree. The MiK\TeX{} equivalent is a user root
registered with \texttt{initexmf}.

\subsubsection*{Route 3 --- Overleaf}

Upload the \file{.sty} and \file{.cls} files into the project alongside the
document, and set the project's compiler to LuaLaTeX in
\textsf{Menu $\rightarrow$ Compiler}. Overleaf resolves files from the project
root, so this is Route 1 with an upload.

\section{A first document}
\label{sec:first}

A complete résumé, in the shape every document in the toolkit follows: load a
class, set the profile once, render the header, then write content.

\begin{latexcode}
\documentclass{careerdossier-resume}

\CDossierSetup{
  name     = {Ada Lovelace},
  headline = {Data Scientist},
  email    = {ada@example.com},
  phone    = {+1 555 555 5555},
  location = {London, UK},
  website  = {example.com},
  linkedin = {ada-lovelace},
  github   = {ada-lovelace}
}

\begin{document}

\MakeCDossierHeader

\CDossierSection{Experience}

\begin{CDossierEntry}[
  title        = {Senior Data Scientist},
  organization = {Example Organization},
  location     = {Toronto, Ontario},
  dates        = {2024--Present}
]
  \begin{CDossierItemize}
    \item Developed a reproducible analytical workflow.
    \item Led the migration of the reporting pipeline.
  \end{CDossierItemize}
\end{CDossierEntry}

\end{document}
\end{latexcode}

Build it with:

\begin{plaincode}
latexmk -lualatex resume.tex
\end{plaincode}

Two passes are required, and latexmk performs them. A résumé and a CV record
each list's length in the auxiliary file so the page-break policy of
\S\ref{sec:pagination} can be applied, and the \texttt{Page N of M} folio needs
the final page count. On a first build from a clean tree the page breaks and the
folio are therefore provisional; the second pass settles both.

The other document types follow the same shape and are shown complete in
\S\ref{sec:examples}.

\section{The document classes}
\label{sec:classes}

\subsection{Loading a class}

The four classes are loaded by name. They live at the top level of the
distribution, so do not write a path:

\begin{latexcode}
\documentclass{careerdossier-resume}
\documentclass{careerdossier-letter}
\documentclass{careerdossier-cv}
\documentclass[type=research]{careerdossier-statement}
\end{latexcode}

\subsection{Supported configuration}

\begin{longtable}{@{}L{0.36\textwidth}L{0.58\textwidth}@{}}
\toprule
\textbf{Setting} & \textbf{Support} \\
\midrule
\endhead
Engine & LuaLaTeX only \\
Language & English \\
Paper & US Letter (default) and opt-in A4 \\
Body size & \opt{fontsize=10pt|11pt|12pt}, per-class default \\
Margins & \opt{margin=normal} (1\,in) or \opt{margin=narrow} (0.5\,in) \\
Body font & Serif (default) and opt-in sans \\
Output medium & \opt{medium=print} (default) or \opt{medium=screen} \\
Entry-metadata de-emphasis & \opt{muted=plain} (default), \opt{italic}, \opt{gray}, \opt{both} \\
Entry-metadata placement & \opt{entrymeta=column} (default) or \opt{inline}; résumé and CV \\
Theme & Monochrome \\
Tagged structure & Opt-in, off by default \\
Résumé class & \file{careerdossier-resume} \\
CV class & \file{careerdossier-cv} \\
Letter class & \file{careerdossier-letter}, industry and academic families \\
Statement class & \file{careerdossier-statement}, default interest type plus six others \\
Bibliography & Optional \file{careerdossier-biblatex} \\
Manual publications & \env{CDossierPublications}, in the CV \\
RTL or bilingual layout & Not supported \\
\bottomrule
\end{longtable}

\subsection{Options and their defaults}

Every option takes a fixed set of values. An unsupported value stops the build
with a class error naming the option, the value supplied, the owning class, and
the accepted values --- so you never have to look the accepted set up here in
order to act on the error. An option is rejected rather than silently ignored.

\begin{longtable}{@{}L{0.16\textwidth}L{0.30\textwidth}L{0.16\textwidth}L{0.30\textwidth}@{}}
\toprule
\textbf{Option} & \textbf{Accepted values} & \textbf{Classes} & \textbf{Default} \\
\midrule
\endhead
\opt{fontsize} & \opt{10pt}, \opt{11pt}, \opt{12pt} & all four & \opt{11pt} résumé; \opt{12pt} others \\
\opt{margin} & \opt{normal}, \opt{narrow} & all four & \opt{narrow} résumé; \opt{normal} others \\
\opt{paper} & \opt{letter}, \opt{a4} & all four & \opt{letter} \\
\opt{bodyfont} & \opt{serif}, \opt{sans} & all four & \opt{serif} \\
\opt{medium} & \opt{print}, \opt{screen} & all four & \opt{print} \\
\opt{muted} & \opt{plain}, \opt{italic}, \opt{gray}, \opt{both} & all four & \opt{plain} \\
\opt{entrymeta} & \opt{column}, \opt{inline} & résumé, CV & \opt{column} \\
\opt{family} & \opt{industry}, \opt{academic} & letter & \opt{industry} \\
\opt{type} & seven values, \S\ref{sec:statement} & statement & \opt{interest} \\
\bottomrule
\end{longtable}

Two settings are fixed and are not options: the language is English and the
theme is monochrome.

\subsubsection{\texttt{fontsize}}
\label{sec:fontsize}

\opt{fontsize} selects one whole-point type scale for the whole document. It is
the only input to type size; there is no per-element size option. The sizes it
resolves to, in points:

\begin{longtable}{@{}lrrr@{}}
\toprule
\textbf{Element} & \opt{10pt} & \opt{11pt} & \opt{12pt} \\
\midrule
\endhead
Name & 19 & 21 & 23 \\
Statement title & 15 & 16 & 18 \\
Headline, subtitle & 12 & 13 & 14 \\
Section heading & 11 & 12 & 13 \\
Entry title, body text, bullets, dates & 10 & 11 & 12 \\
Contact line & 9 & 10 & 11 \\
Running header, folio & 8 & 9 & 10 \\
\bottomrule
\end{longtable}

Structural vertical spacing scales with the same option, as a fixed fraction of
the body baseline. \S\ref{sec:tokens} lists the tokens that name each gap.

\opt{density=compact|standard} was removed in \texttt{v0.6.0} and has no
replacement.

\subsubsection{\texttt{margin}}

\opt{normal} means one-inch margins and \opt{narrow} half-inch margins.
\opt{margin} and \opt{fontsize} together decide line length. Measured in TeX
Gyre Termes on US Letter, full lines of running prose hold roughly:

\begin{longtable}{@{}lll@{}}
\toprule
\opt{fontsize} & \opt{margin} & \textbf{Characters per line} \\
\midrule
\endhead
\opt{10pt} & \opt{normal} & 110--120 \\
\opt{10pt} & \opt{narrow} & 130--140 \\
\opt{11pt} & \opt{normal} & 102--113 \\
\opt{11pt} & \opt{narrow} & 118--127 \\
\opt{12pt} & \opt{normal} & 93--101 \\
\opt{12pt} & \opt{narrow} & 108--117 \\
\bottomrule
\end{longtable}

The prose classes default to \opt{12pt} at \opt{margin=normal} because that is
the only combination whose measure lands near the conventional 45--90 character
guidance at a one-inch margin.

\paragraph{Known limitation.} The résumé's default \opt{11pt} with
\opt{margin=narrow} produces the longest measure in the project, roughly
118--127 characters per line. It is chosen for one-page capacity and is fine for
entry lines and short bullets, which are split across the measure and never
reach its full width. It is visibly long in a full-width summary paragraph. If a
résumé opens with a dense prose summary, prefer \opt{margin=normal},
\opt{fontsize=12pt}, or a shorter summary.

The classes load \textsf{geometry}. An advanced user may call \cmd{geometry}
after \cmd{documentclass} to replace the preset. Do not load \textsf{geometry}
again with package options: that can produce an option clash. A raw override
bypasses the two tested presets, so its pagination and visual result are the
author's responsibility.

\subsubsection{\texttt{paper}}

\opt{a4} selects an ISO A4 media box while retaining the selected margin preset,
font size, spacing, and page-furniture design. Because A4 is slightly narrower
and taller than US Letter, line and page breaks may change.

\subsubsection{\texttt{bodyfont}}

\opt{serif}, the default, preserves the TeX Gyre Termes body with TeX Gyre Heros
headings. \opt{sans} selects TeX Gyre Heros for the document body; headings
remain TeX Gyre Heros. The option changes no size, semantic role, spacing,
geometry, or page furniture. Arbitrary font names and per-role font selection
are not supported.

\subsubsection{\texttt{medium}}

The same dossier is read in two contexts with different needs. On screen the PDF
viewer already shows page position, so a folio is redundant chrome; on paper a
loose page has no such affordance, and the folio is what keeps a multi-page
dossier in order after it is put down.

\opt{print}, the default, emits the page furniture described in
\S\ref{sec:furniture}. \opt{screen} emits no running header and no folio on any
page.

\opt{medium} also decides whether a link carries a visible decoration:

\begin{longtable}{@{}L{0.20\textwidth}L{0.70\textwidth}@{}}
\toprule
\textbf{Value} & \textbf{Author-written \cmd{href} anchor text} \\
\midrule
\endhead
\opt{print} & undecorated --- reads as ordinary body text \\
\opt{screen} & underlined \\
\bottomrule
\end{longtable}

Links are drawn in body black with no border under both values, which is the
print-safe treatment. Where a link's visible text is the address --- the contact
line, \cmd{CDossierLink}, an ORCID iD, a bibliography DOI --- the absence of a
rule costs a reader nothing, because the text announces itself. Where you supply
your own anchor text there is otherwise no cue of any kind:

\begin{latexcode}
Shipped the analytical engine, documented in a
\href{https://example.org/work}{public write-up}.
\end{latexcode}

Under \opt{medium=screen} \emph{public write-up} is underlined; under
\opt{medium=print} it is not. A rule under a contact line's e-mail and website
but not under its phone or location would read as emphasis rather than as
linking, which is why address-as-text links are never decorated.
\cmd{CDossierPlainLinks} suppresses the decoration for a region.

The link \emph{annotation} is emitted under both media. It is invisible on paper
and costs nothing there, and stripping it would hand dead links to anyone who
opens a \opt{medium=print} file on screen.

The option changes only whether furniture and decoration are emitted. Page
geometry is unchanged, so switching \opt{medium} cannot reflow a document, and
extracted text is identical under both values.

\subsubsection{\texttt{muted}}

\opt{muted} decides how de-emphasised runs are rendered: an entry's dates and
location in the résumé and CV, and the statement's application-context line.
Note the spelling \opt{gray}, not \emph{grey}.

\begin{longtable}{@{}L{0.28\textwidth}L{0.24\textwidth}L{0.40\textwidth}@{}}
\toprule
\textbf{Value} & \textbf{Shape} & \textbf{Colour} \\
\midrule
\endhead
\opt{plain} (default) & upright & text token (black) \\
\opt{italic} & italic & text token (black) \\
\opt{gray} & upright & muted token (gray 0.30) \\
\opt{both} & italic & muted token (gray 0.30) \\
\bottomrule
\end{longtable}

The muted token measures 8.52:1 against white under the WCAG 2.1
relative-luminance formula, well above the 4.5:1 normal-text floor.

Which styling to want is a genuine trade-off, which is why it is an option.
Italic at small sizes is harder to read for low-vision and dyslexic readers than
a high-contrast gray, and this metadata is scanned rather than read; but shape,
unlike a gray level, survives a fax, a photocopy, and a 1-bit print.
\opt{plain} is the default because both halves are real, and a reader affected
by both would otherwise have nothing to select. Under \opt{plain} the dates and
location are distinguished by position and content alone, which is the design's
standing claim under every value.

All four values are visual only: no extractor sees a difference, and the reading
order is identical under each. Colour is never the only carrier of meaning.

\subsubsection{\texttt{entrymeta}}

Accepted by the résumé and CV; the letter and statement classes do not accept it
because neither has entry headings.

\opt{entrymeta} decides where an entry's dates and location sit. Under the
default \opt{column} they are set flush right, opposite the title and
organization:

\begin{plaincode}
Senior Engineer                                          2024-Present
Example Labs                                              Toronto, ON
  * First achievement in source order.
\end{plaincode}

Under \opt{inline} they follow the title and organization on the same line,
joined by \cmd{CDossierEntryMetaSeparator}:

\begin{plaincode}
Senior Engineer | 2024-Present
Example Labs | Toronto, ON
  * First achievement in source order.
\end{plaincode}

Nothing else changes under either value: the cells keep their order, their
semantic roles, and the page-break penalty that holds the two heading lines
together, and a missing \key{organization}, \key{location}, or \key{dates}
leaves no stray separator under \opt{inline} just as it leaves no empty column
under \opt{column}.

\paragraph{What the option is for.} \cmd{CDossierRecordListEdgeAboveSkip}
carries a lower bound of \texttt{0.25} that is an extraction constraint rather
than a design preference (\S\ref{sec:itemize}). The bound exists because the
flush-right column makes each heading row a two-column region for a geometric
extractor, and only enough vertical separation below the heading keeps the dates
from being read after the bullets. \opt{inline} has no column, so no such region
arises and the bound does not apply to a document that selects it. This is the
only supported way to set a list edge below \texttt{0.25} without breaking
reading order:

\begin{latexcode}
\documentclass[entrymeta=inline]{careerdossier-resume}
\ExplSyntaxOn
\skip_set:Nn \CDossierRecordListEdgeAboveSkip
  { \fp_to_dim:n { 0.125 * \dim_to_fp:n { \CDossierBodyLeading } } }
\ExplSyntaxOff
\end{latexcode}

Under \opt{column} that override still breaks reading order, and nothing detects
it for you.

\opt{inline} is not the default because it is a visible design change --- it
discards the column the résumé's layout is built around --- and because making it
the default would reflow every existing document to buy a constraint only some
documents care about. Design, a clean text layer, and a sub-\texttt{0.25} list
edge are mutually exclusive; any two are available.

Both values are equivalent on the tagged path: the separator is emitted as a
layout artifact, exactly as the contact line's \texttt{|} is, so a consumer
reading the structure tree receives the same logical text under either value and
assistive technology is not made to announce a vertical bar.

\paragraph{\cmd{CDossierEntryMetaSeparator}} is the mark \opt{entrymeta=inline}
places between two cells, with its surrounding space. The default is
\verb|~|\verb+|+\verb|~|, the same mark the contact line uses, so a document has
one field separator rather than two. It is emitted only between two cells that
are both present, so redefining it cannot introduce a stray separator. Two
properties of the default are worth preserving in a replacement:

\begin{itemize}
  \item The mark is wrapped as a \textbf{layout artifact}, which is what keeps
        it out of the structure tree on the tagged path. A replacement set as
        ordinary text renders identically and will also be announced by a screen
        reader.
  \item The \textbf{leading} space is ordinary content glue and the
        \textbf{trailing} one sits inside the artifact. An artifact interrupts
        the enclosing run, so only the leading space survives as the word
        boundary between the two cells; a second content space would show up as
        a doubled gap in the extracted text.
\end{itemize}

Keep the mark out of the muted role as well: the cell following it is set in
\cmd{CDossierMutedStyle}, and an italic separator reads as a glyph belonging to
the metadata rather than as a boundary between fields. The default wraps the
mark in \cmd{CDossierBodyStyle}, which resets shape as well as family.

\subsubsection{\texttt{family} (letter)}

\opt{family} accepts \opt{industry} and \opt{academic}, defaulting to
\opt{industry}. It is label- and metadata-only: it selects the PDF document type
while both values use the same geometry, calibrated type scale, prose rhythm,
letterhead structure, and shared \texttt{Cover Letter} continuation label.
Switching between the two does not select a denser or roomier layout, and does
not introduce new recipient keys.

\subsection{Loading a package directly}

Three of the packages also accept options when loaded with \cmd{usepackage}, and
report an unsupported value the same way a class does:

\begin{longtable}{@{}L{0.30\textwidth}L{0.30\textwidth}L{0.32\textwidth}@{}}
\toprule
\textbf{Package} & \textbf{Options} & \textbf{Package default} \\
\midrule
\endhead
\file{careerdossier-tokens} & \opt{fontsize}, \opt{margin} & \opt{12pt}, \opt{normal} \\
\file{careerdossier-typography} & \opt{bodyfont} & \opt{serif} \\
\file{careerdossier-components} & \opt{medium}, \opt{muted}, \opt{entrymeta} & \opt{print}, \opt{plain}, \opt{column} \\
\bottomrule
\end{longtable}

\textbf{A package default is not a class default.} Loading
\file{careerdossier-tokens} directly gives \opt{12pt} and \opt{normal}, where
the résumé class resolves to \opt{11pt} and \opt{narrow} --- because the class
passes its own values down before it loads the package. Set the option
explicitly when loading a package directly; do not assume it inherits a class's
choice.

\section{Profile metadata}
\label{sec:profile}

\subsection{\cmd{CDossierSetup}}

One profile serves every document. Set it once, in the preamble:

\begin{latexcode}
\CDossierSetup{
  name     = {Ada Lovelace},
  headline = {Data Scientist},
  email    = {ada@example.com},
  phone    = {+1 555 555 5555},
  location = {London, UK},
  website  = {example.com},
  linkedin = {ada-lovelace},
  github   = {ada-lovelace},
  scholar  = {kukA0LcAAAAJ},
  orcid    = {0000-0002-1825-0097}
}
\end{latexcode}

In practice the profile lives in its own file, included by every document built
from it, so that an address is corrected in one place.

\subsection{Profile keys}

\begin{longtable}{@{}llL{0.52\textwidth}@{}}
\toprule
\textbf{Key} & \textbf{Required} & \textbf{Purpose} \\
\midrule
\endhead
\key{name} & Yes & Person's display name \\
\key{headline} & No & Professional title or short descriptor \\
\key{email} & No & E-mail address \\
\key{phone} & No & Telephone number \\
\key{location} & No & City, region, or country \\
\key{website} & No & Personal or professional website \\
\key{linkedin} & No & LinkedIn handle, profile path, or URL \\
\key{github} & No & GitHub handle, profile path, or URL \\
\key{scholar} & No & Google Scholar identifier, profile path, or URL \\
\key{orcid} & No & ORCID identifier or profile URL \\
\key{affiliation} & Conditional & Current institution, organization, studio, or independent-practice description \\
\bottomrule
\end{longtable}

Whitespace-only values count as missing.

\key{name} is the only key every class requires. A class may require a further
key for a particular document, and the statement class does: \key{email} for
every statement, \key{affiliation} for \opt{type=research}, and \key{website}
for \opt{type=artist}. \emph{Conditional} above means exactly that --- the key is
optional to the profile and required by some documents that use it.

A shared profile is expected to carry fields for several document types, so a
class that does not display a field does not warn about it.

\subsection{Accepted forms for the web-profile keys}
\label{sec:webkeys}

\key{linkedin}, \key{github}, and \key{scholar} each name a service with a
documented canonical address, so each accepts a bare identifier as well as an
address. A bare identifier is expanded to the canonical address --- for the
displayed text and the link target alike, so the address on the page is the
address the link goes to:

\begin{longtable}{@{}L{0.14\textwidth}L{0.24\textwidth}L{0.54\textwidth}@{}}
\toprule
\textbf{Key} & \textbf{Written} & \textbf{Displayed} \\
\midrule
\endhead
\key{linkedin} & \texttt{ada-lovelace} & \addr{linkedin.com/in/ada-lovelace} \\
\key{github} & \texttt{ada-lovelace} & \texttt{github.com/ada-lovelace} \\
\key{scholar} & \texttt{kukA0LcAAAAJ} & \addr{scholar.google.com/citations?user=kukA0LcAAAAJ} \\
\bottomrule
\end{longtable}

A value is read as a bare identifier when it contains neither \texttt{/} nor
\texttt{.}\ --- the two characters every address form carries, and that none of
these three services permits in an identifier.

Any other value is used exactly as written, with only the \texttt{https://}
scheme supplied when it carries none. A \texttt{www.}\ host, a trailing slash,
and query parameters are all preserved.

\key{website} has no canonical host by definition and is never expanded: a value
with no \texttt{/} or \texttt{.}\ is displayed and linked exactly as written.
\key{orcid} keeps its own normalization and its \texttt{ORCID:} label
(\S\ref{sec:orcid}).

\subsubsection{Forms that work}

Give either the identifier alone or a complete address. The forms in between are
what go wrong.

\begin{longtable}{@{}L{0.13\textwidth}L{0.40\textwidth}L{0.39\textwidth}@{}}
\toprule
\textbf{Key} & \textbf{Write} & \textbf{Displayed and linked as} \\
\midrule
\endhead
\key{linkedin} & \texttt{ada-lovelace} & \addr{linkedin.com/in/ada-lovelace} \\
\key{linkedin} & \addr{linkedin.com/in/ada-lovelace} & unchanged \\
\key{linkedin} & \addr{www.linkedin.com/in/ada-lovelace} & unchanged, \texttt{www.}\ kept \\
\key{linkedin} & \addr{https://www.linkedin.com/in/ada-lovelace/} & unchanged, scheme and trailing slash kept \\
\key{github} & \texttt{ada-lovelace} & \texttt{github.com/ada-lovelace} \\
\key{github} & \texttt{github.com/ada-lovelace} & unchanged \\
\key{scholar} & \texttt{kukA0LcAAAAJ} & \addr{scholar.google.com/citations?user=kukA0LcAAAAJ} \\
\key{scholar} & \addr{scholar.google.com/citations?user=kukA0LcAAAAJ} & unchanged \\
\bottomrule
\end{longtable}

A Google Scholar identifier containing \texttt{\_} needs no escape.

\subsubsection{Forms to avoid}

These build successfully and produce a broken link, which is the case worth
knowing about: the page looks correct and nothing in the log says otherwise.
Bare identifiers are not validated, because nothing in \LaTeX{} can dereference
an address.

\begin{longtable}{@{}L{0.30\textwidth}L{0.26\textwidth}L{0.36\textwidth}@{}}
\toprule
\textbf{Written} & \textbf{You get} & \textbf{Why} \\
\midrule
\endhead
\texttt{linkedin = \{in/ada-lovelace\}} & \addr{https://in/ada-lovelace} & half a path: the \texttt{/} makes it read as an address, so it is left alone \\
\texttt{linkedin = \{ada.lovelace\}} & \addr{https://ada.lovelace} & the \texttt{.}\ makes it read as a host, and no LinkedIn slug contains one \\
\texttt{linkedin = \{Ada Lovelace\}} & \addr{linkedin.com/in/Ada Lovelace} & a display name, not a handle; note the space inside the address \\
\texttt{github = \{@ada-lovelace\}} & \addr{github.com/@ada-lovelace} & \texttt{@} is a social-media convention, not part of a GitHub address \\
\texttt{github = \{ada-lovelace/some-repo\}} & \addr{https://ada-lovelace/some-repo} & a repository, not a profile --- put repository links in body text with \cmd{CDossierLink} \\
\texttt{scholar = \{user=kukA0LcAAAAJ\}} & the prefix added on top of it & the query rather than the identifier \\
\bottomrule
\end{longtable}

\subsubsection{Characters that stop the build}
\label{sec:escaping}

A field value is read as ordinary text, not as a verbatim argument, so one
character is consumed by \TeX{} before the link machinery ever sees it and must
be escaped. This applies to every profile field, not only the three above.

\begin{longtable}{@{}L{0.26\textwidth}L{0.16\textwidth}L{0.50\textwidth}@{}}
\toprule
\textbf{Written} & \textbf{Result} & \textbf{Write instead} \\
\midrule
\endhead
\texttt{example.com/a\%b} & build fails & \texttt{example.com/a\textbackslash\%b} \\
\texttt{example.com/a\%25b} & build fails & \texttt{example.com/a\textbackslash\%25b} \\
\bottomrule
\end{longtable}

Percent-encoding does not help; the \texttt{\%} is still a comment character.
This cannot be fixed at the package level, and the reason is worth stating once:
\TeX's lexer discards \texttt{\%} and the rest of the line while
\cmd{CDossierSetup} is still reading its argument, so the tokens never reach any
code this package could write. \verb|\%| is the only possible answer.

Nothing else needs escaping. \texttt{\#}, \texttt{\&}, \texttt{\_},
\texttt{\textasciitilde}, \texttt{\$}, and \texttt{\textasciicircum} all build
and render correctly, so a pasted Google Scholar address carrying
\texttt{\&hl=en} is fine, and so is a link carrying \texttt{\#installation}.

\textbf{Escaping them anyway is harmless}, which makes the rule safe to apply
defensively. Written either way, \texttt{a\&b} and \texttt{a\textbackslash\&b}
produce a byte-identical link action in the PDF. So the whole rule is
\emph{escape} \texttt{\%}; \emph{escaping anything else changes nothing}.
Over-escaping cannot silently break a link.

In practice this affects \key{website} and \cmd{CDossierLink} far more than the
three keys above, since a LinkedIn or GitHub handle and a Scholar identifier are
restricted to letters, digits, hyphens, and underscores.

\subsection{ORCID}
\label{sec:orcid}

An ORCID iD is displayed as ordinary text with a descriptive \texttt{ORCID:}
label and a web link. A bare identifier is normalized to an
\texttt{https://orcid.org/} target; a complete URL is used as supplied. ORCID is
the one web-profile key whose displayed text stays the bare identifier: its
label supplies the meaning the other three take from their domain.

The value is checked against the shape of an ORCID iD --- sixteen digits in four
hyphen-separated groups whose final character is a digit or \texttt{X} --- and a
value matching neither that shape nor an ORCID address raises a warning naming
the value and the expected form:

\begin{longtable}{@{}L{0.46\textwidth}L{0.46\textwidth}@{}}
\toprule
\textbf{Written} & \textbf{Result} \\
\midrule
\endhead
\texttt{0000-0002-1825-0097} & accepted \\
\texttt{0000-0002-1825-009X} & accepted --- \texttt{X} is the documented checksum character \\
\texttt{orcid.org/0000-0002-1825-0097} & accepted \\
\addr{https://orcid.org/0000-0002-1825-0097} & accepted \\
\texttt{0000-0002-1825-009} & warns --- too short \\
\texttt{0000-0002-1825-009X7} & warns --- too long \\
\addr{https://example.org/ada-lovelace} & warns --- not an ORCID address \\
\addr{scholar.google.com/citations?user=example} & warns --- a value in the wrong key \\
\bottomrule
\end{longtable}

An optional \texttt{www.}, a lowercase \texttt{x} checksum character, and one
trailing slash are accepted without a warning: they resolve, and warning about a
working value would only teach authors to ignore the warning that matters.

The check never stops the build, and nothing about the stored value, the link
target, or the label depends on its outcome --- an author holding an identifier
this check does not anticipate still gets exactly the document they wrote. It is
a shape check only: it does not verify the ISO 7064 checksum, and it cannot tell
whether an iD is registered, since \LaTeX{} has no network access.

\subsection{Contact-field labels}
\label{sec:contactlabels}

\begin{latexcode}
\CDossierSetup{ contact-labels = true }   % default: false
\end{latexcode}

\key{contact-labels} is an option key, not a profile field: it holds no value to
print and cannot be read back with \cmd{CDossierPrintField}. When enabled, the
contact line prefixes a short identifying label to the fields whose nature the
value itself does not convey:

\begin{longtable}{@{}L{0.22\textwidth}L{0.70\textwidth}@{}}
\toprule
\textbf{Field} & \textbf{Rendered as} \\
\midrule
\endhead
\key{email} & \texttt{Email: ada@example.com} \\
\key{phone} & \texttt{Phone: +1 555 0100} \\
\key{website} & \texttt{Website: example.test/resume} \\
\bottomrule
\end{longtable}

The remaining fields stay unlabelled by design: \key{linkedin}, \key{github},
and \key{scholar} display beginning with their service's domain --- including
when you supplied a bare identifier, since that is expanded --- \key{orcid}
always carries its own label, and a \key{location} value is a place name. Labels
are fixed English strings; the project is English-only.

\paragraph{Why it exists.} In the default rendering a screen reader announces
the e-mail and the website as links, so their nature is conveyed --- but the
phone number is announced as bare digits with nothing indicating what it is.
Sighted readers infer all three from format and position; that inference is
exactly what a screen-reader user does not get. A visible text label is the one
mechanism that works in every consumer --- screen readers, plain text
extraction, and ATS parsers --- including the default untagged output, which is
why it is the primary fix rather than a tag-level attribute.

Behavioural guarantees:

\begin{itemize}
  \item the default rendering is unchanged, and the feature is strictly opt-in;
  \item the key applies to every class's contact line, including
        document-specific subsets such as the statement classes';
  \item an absent field leaves no orphan label and no stray separator;
  \item labels are content, not layout artifacts, so they reach the structure
        tree and the extracted text; separators remain artifacts;
  \item \key{contact-labels} alone means \texttt{true}, and the value must be a
        boolean.
\end{itemize}

\subsection{Contact-line wrapping}

Contact fields form a centred block whose rows are packed greedily in the fixed
field order. Complete fields are the only permitted units: e-mail addresses,
phone numbers, locations, websites, LinkedIn, GitHub, Scholar, and ORCID items
remain intact on one visual line. A separator is inserted only after two
adjacent items are known to fit on the same row, so no visual line begins or
ends with one.

The component does not truncate values or reduce the selected font size. Each
complete item is measured against the available width, and an item that cannot
fit on an otherwise empty line raises an error naming the field. Shorten its
displayed value, select \opt{margin=narrow}, or choose a smaller supported
\opt{fontsize}.

\subsection{Required-field validation}

\key{name} is required before rendering a dossier header or letterhead. A
missing name produces an error that names the missing key, identifies the
command that required it, and shows a minimal correction. Optional fields never
produce empty separators or blank lines.

\section{Letter metadata}
\label{sec:letter}

\begin{latexcode}
\CDossierLetterSetup{
  date                   = {\today},
  recipient-name         = {Hiring Manager},
  recipient-title        = {Director of Analytics},
  recipient-organization = {Example Organization},
  recipient-address      = {123 Example Street\\Toronto, Ontario},
  subject                = {Application for the Data Scientist Position},
  salutation             = {Dear Hiring Manager,},
  closing                = {Sincerely,}
}
\end{latexcode}

\begin{longtable}{@{}llL{0.46\textwidth}@{}}
\toprule
\textbf{Key} & \textbf{Required} & \textbf{Default or behaviour} \\
\midrule
\endhead
\key{date} & No & \cmd{today} \\
\key{recipient-name} & No & Omitted when empty \\
\key{recipient-title} & No & Omitted when empty \\
\key{recipient-organization} & No & Omitted when empty \\
\key{recipient-address} & No & Omitted when empty; may contain \texttt{\textbackslash\textbackslash} \\
\key{subject} & No & Entire subject line omitted when empty \\
\key{salutation} & No & \texttt{Dear Hiring Manager,} \\
\key{closing} & No & \texttt{Sincerely,} \\
\bottomrule
\end{longtable}

The recipient block collapses cleanly when one or more optional fields are
absent: an absent block or subject omits both the content and its following
block skip.

\section{Statements}
\label{sec:statement}

All statement documents use one class with an optional \opt{type} option and the
shared settings. When \opt{type} is omitted the class selects \opt{interest},
which requires no profile fields beyond \key{name} and \key{email}. Supplying
\opt{type} without a value, or with an unsupported value, produces an actionable
class error.

\begin{longtable}{@{}lL{0.42\textwidth}l@{}}
\toprule
\opt{type} & \textbf{Default title} & \textbf{Default running title} \\
\midrule
\endhead
\opt{interest} (default) & Statement of Interest & Statement of Interest \\
\opt{research} & Research Statement & Research Statement \\
\opt{teaching} & Teaching Statement & Teaching Statement \\
\opt{teaching-philosophy} & Statement of Teaching Philosophy & Teaching Philosophy \\
\opt{diversity} & Statement of Contributions to Equity, Diversity, Inclusion, and Accessibility & EDIA Statement \\
\opt{artist} & Artist Statement & Artist Statement \\
\opt{purpose} & Statement of Purpose & Statement of Purpose \\
\bottomrule
\end{longtable}

The type selects a title, continuation-page identification, and required-field
contract. It does not generate or enforce content sections.

\subsection{Statement metadata}

Document-specific values use a separate setup command, because they describe one
statement rather than a reusable profile:

\begin{latexcode}
\CDossierStatementSetup{
  title               = {Research Statement},
  running-title       = {Research Statement},
  subtitle            = {Reliable scientific computing},
  application-context = {Application for Assistant Professor of Computational Science},
  application-id      = {APP-2026-0042}
}
\end{latexcode}

\begin{longtable}{@{}llL{0.50\textwidth}@{}}
\toprule
\textbf{Key} & \textbf{Required} & \textbf{Default or behaviour} \\
\midrule
\endhead
\key{title} & No & Default selected by the active \opt{type}; a nonblank value overrides it \\
\key{running-title} & No & Short default selected by the active \opt{type}; overrides independently of \key{title} \\
\key{subtitle} & No & One short line beneath the title; omitted when blank \\
\key{application-context} & No & Separate contextual line; omitted when blank \\
\key{application-id} & No & Rendered as labelled text with the application context; omitted when blank \\
\bottomrule
\end{longtable}

Whitespace-only values count as absent. Later setup calls overwrite earlier
values.

Current affiliation is reusable identity data rather than application-specific
state, so it is a shared profile key (\key{affiliation}) rather than a statement
key.

\subsection{Required fields and displayed contacts}

Every statement requires \key{name} and \key{email}. The header validates the
following additional type-specific requirements and renders only the listed
optional contacts:

\begin{longtable}{@{}L{0.26\textwidth}L{0.24\textwidth}L{0.42\textwidth}@{}}
\toprule
\textbf{Type} & \textbf{Additional required field} & \textbf{Optional displayed fields} \\
\midrule
\endhead
\opt{interest} & None & \key{phone}, \key{website}, \key{affiliation} \\
\opt{research} & \key{affiliation} & \key{phone}, \key{website}, \key{scholar}, \key{orcid} \\
\opt{teaching} & None & \key{phone}, \key{website}, \key{affiliation} \\
\opt{teaching-philosophy} & None & \key{phone}, \key{website}, \key{affiliation} \\
\opt{diversity} & None & \key{phone}, \key{website}, \key{affiliation} \\
\opt{artist} & \key{website} & \key{phone}, \key{affiliation} \\
\opt{purpose} & None & \key{phone}, \key{website}, \key{affiliation} \\
\bottomrule
\end{longtable}

For \opt{artist}, \key{website} may identify a personal site, portfolio, or
comparable public presence. Shared \key{headline}, \key{location},
\key{linkedin}, and \key{github} values remain available to other documents but
are not displayed by the statement header, and their presence does not warn.

\subsection{First-page identity order}

\cmd{MakeCDossierStatementHeader} validates the active type and emits the
present items in this fixed logical order:

\begin{enumerate}
  \item profile \key{name};
  \item the selected statement title;
  \item optional one-line \key{subtitle};
  \item optional or required profile \key{affiliation};
  \item optional \key{application-context}, followed by a labelled
        \key{application-id} when both are present;
  \item required \key{email}, followed by the active type's present optional
        contacts.
\end{enumerate}

Application context is not a second subtitle and is not mixed into the contact
list. An absent value leaves neither a blank line nor a gap. The page-one title
drives the PDF document metadata; the running title exists only for page
furniture.

\subsection{Author content and headings}

The class introduces no command for research aims, teaching themes, EDI
commitments, artistic methods, or statement-of-purpose paragraphs. Authors write
ordinary prose and may add headings where the content benefits. Two unnumbered
heading levels are available, spelled exactly as in the résumé and CV:

\begin{latexcode}
\CDossierSection{Research Vision and Significance}
\CDossierSubsection{A Narrower Theme}
\end{latexcode}

The rhythm is this class's own. A statement separates its paragraphs by a
visible \cmd{CDossierProseParSkip}, so its heading gaps are calibrated against
that and use the prose tokens rather than the record classes':

\begin{longtable}{@{}L{0.24\textwidth}L{0.68\textwidth}@{}}
\toprule
\textbf{Property} & \textbf{Statement heading} \\
\midrule
\endhead
Section size & \cmd{CDossierSizeSection} --- the section-heading step of the
  type scale in \S\ref{sec:fontsize} \\
Subsection size & \cmd{CDossierSizeBody} --- the body step of the same scale \\
Typeface & \cmd{CDossierSectionStyle}, in both \opt{bodyfont} modes \\
Space above & \cmd{CDossierProseSectionAboveSkip} / \cmd{CDossierProseSubsectionAboveSkip} \\
Space below & \cmd{CDossierProseSectionBelowSkip} / \cmd{CDossierProseSubsectionBelowSkip} \\
Decorative rule & None \\
\bottomrule
\end{longtable}

A statement heading carries no rule: the full-width rule under
\cmd{CDossierSection} in the résumé and CV is entry-structured section
furniture, and a prose document reads better without it.

Numbered sectioning and levels below \cmd{subsection} are inherited from
\textsf{article} unchanged and are outside the calibrated design.

\section{Content commands and environments}
\label{sec:content}

\subsection{\cmd{MakeCDossierHeader}}

Renders the résumé or CV identity block from the shared profile. It validates
\key{name}; sets the name, optional headline, and contact line on the calibrated
type scale; renders \key{headline} only when present and leaves no gap behind
when it is absent; renders the available contact fields, inserting separators
only between rendered fields that remain on the same visual line; creates links
for the supported contact fields; prefixes the text labels of
\S\ref{sec:contactlabels} when enabled; and adds no page numbers.

\subsection{\cmd{MakeCDossierLetterhead} and \cmd{MakeCDossierClosing}}

\cmd{MakeCDossierLetterhead} renders the cover-letter opening material in order:
the shared sender identity, the date, the recipient block, the optional subject,
and the salutation. An absent recipient block or subject omits both the content
and its following block skip.

\cmd{MakeCDossierClosing} renders the configured closing, suitable signature
space, and the profile \key{name}, and validates that \key{name} exists.

The three are a single unit for pagination: the closing text and the signature
name always land on the same page, whichever page that turns out to be. When
they do not both fit under the body, the page breaks \emph{above} the closing and
the whole block opens the next page. This is a contract of the command rather
than a token --- the boundary between the two is forbidden outright, because
under \cmd{raggedbottom} a page-break penalty here is a boolean rather than a
dial (\S\ref{sec:penalties}), and a keep-together that is not infinite is not a
keep. It is a separate guarantee from \cmd{clubpenalty} and \cmd{widowpenalty},
which govern the first and last line \emph{of one paragraph} and cannot express
it.

\subsection{The header stack}

\begin{latexcode}
\CDossierHeaderBegin
\CDossierHeaderLine{<content>}
\CDossierHeaderLine{<content>}
\CDossierHeaderEnd
\end{latexcode}

Composes a centred header stack line by line. \cmd{MakeCDossierHeader} and
\cmd{MakeCDossierStatementHeader} render fixed line lists and are what a
document normally uses; this triple is the interface beneath them, for a header
whose lines a document has to choose itself.

A caller states only which lines are present, in reading order. Every gap is
owned by the components package, and the triple adds no spacing option: no gap
above the first line, \cmd{CDossierSharedHeaderNameGapSkip} below the first,
\cmd{CDossierSharedHeaderMetaGapSkip} below every later one, and the rendering
class's own header-below token as the boundary under the stack. Position decides
which token guards a boundary, never presence, so an omitted line leaves neither
a blank line nor a gap behind. Emit a line conditionally by calling
\cmd{CDossierHeaderLine} only when you have content for it.

\begin{itemize}
  \item \cmd{CDossierHeaderBegin} starts a stack and discards any lines left
        over from an abandoned one.
  \item \cmd{CDossierHeaderLine} stores its argument unexpanded and typesets
        nothing yet, so a line may carry structure --- a heading, a link, a
        tagged element --- and keeps it.
  \item \cmd{CDossierHeaderEnd} renders the collected lines, each as its own
        centred paragraph, and emits the boundary below the stack.
  \item A stack with no lines renders nothing, including no boundary below it.
  \item No validation is performed: the wrapper commands keep their own.
\end{itemize}

A worked example --- a résumé wanting \key{location} on a line of its own
\emph{between} the headline and the contact line, which is a position no
append-only hook can reach. It is a complete document, and
\file{tests/smoke/components-header-stack-doc.tex} compiles this exact text on
every run, so the example on this page is one the toolkit is known to accept:

\begin{latexcode}
\documentclass{careerdossier-resume}

\CDossierSetup{
  name     = {Ada Lovelace},
  headline = {Data Scientist},
  location = {London, UK},
  email    = {ada@example.com},
}

\begin{document}

\CDossierHeaderBegin
\CDossierHeaderLine{
  \CDossierNameStyle \CDossierSizeName
  \CDossierPrintField{name}
}
\CDossierHeaderLine{
  \CDossierHeadlineStyle \CDossierSizeHeadline
  \CDossierPrintField{headline}
}
\CDossierIfFieldTF{location}
  {
    \CDossierHeaderLine{
      \CDossierMutedStyle \CDossierSizeSmall
      \CDossierPrintField{location}
    }
  }
  { }
\CDossierHeaderLine{
  \CDossierBodyStyle \CDossierSizeSmall
  \CDossierPrintField{email}
}
\CDossierHeaderEnd

\CDossierSection{Experience}

\end{document}
\end{latexcode}

The guard wraps the \cmd{CDossierHeaderLine} call rather than appearing inside
it: that is the whole of ``conditional''. Drop the \key{location} key and the
header renders three lines with no gap left behind.

\subsection{\cmd{CDossierSection}}

\begin{latexcode}
\CDossierSection{Experience}
\end{latexcode}

Creates a résumé or CV section heading followed by a full-width decorative rule.
Class-controlled spacing places the rule closer to its heading than to the
content below.

The rule sits \cmd{CDossierRecordSectionRuleGapSkip} below the heading's
\emph{baseline}, not below the bottom of its line box, so its height does not
change when the heading happens to contain a descender. The token consequently
has a lower bound: it must exceed the heading's depth, or the rule would cross
descender ink.

The gap between the rule and the section's first content is the larger of
\cmd{CDossierRecordSectionBelowSkip} and whatever leading space the following
block contributes --- never their sum. A section that opens with an entry, a
bullet list, or an ordinary paragraph therefore yields one predictable gap
rather than three different ones.

The argument is user-visible text; the command does not translate it. The
statement class defines \cmd{CDossierSection} too, without a rule and on the
prose heading rhythm.

\subsection{\cmd{CDossierSubsection}}

\begin{latexcode}
\CDossierSection{Experience}
\CDossierSubsection{Industry}
\end{latexcode}

A second-level heading inside a résumé or CV section, for natural groupings that
are not themselves sections --- \emph{Experience} split into industry and
academic, \emph{Publications} split into journal, conference, and preprint.
Without it such a group has to be promoted to a full ruled section, which
flattens the hierarchy it was meant to show.

The heading carries no rule and no size of its own: it is
\cmd{CDossierSectionStyle} at \cmd{CDossierSizeBody}, so the level is marked by
weight, face, and spacing. The rule belongs to the section.

The gap below is \cmd{CDossierRecordSubsectionBelowSkip}, and it is the same gap
whether the next block is a run of entries or a publication list. The gap above
is \cmd{CDossierRecordSubsectionAboveSkip}, with one deliberate exception: a
section that opens \emph{directly} with a subsection does not take that gap, and
the boundary is \cmd{CDossierRecordSectionBelowSkip} alone. Anything typeset in
between ends the exception.

\subsection{\env{CDossierEntry}}

\begin{latexcode}
\begin{CDossierEntry}[
  title        = {Data Scientist},
  organization = {Example Organization},
  location     = {Toronto, Ontario},
  dates        = {2024--Present}
]
  Entry content.
\end{CDossierEntry}
\end{latexcode}

\begin{longtable}{@{}llL{0.50\textwidth}@{}}
\toprule
\textbf{Key} & \textbf{Required} & \textbf{Purpose} \\
\midrule
\endhead
\key{title} & Yes & Position, degree, award, or entry title \\
\key{organization} & No & Employer, institution, or organization \\
\key{location} & No & City, region, country, or remote status \\
\key{dates} & No & Date or date range \\
\bottomrule
\end{longtable}

Named keys are preferred over positional arguments because they are
self-documenting and can be extended without changing the meaning of existing
arguments. Missing optional keys are omitted without leaving visible punctuation
or spacing artifacts. The environment controls the entry heading and local
spacing but does not force bullet content.

An entry body may be empty. A body-less entry --- the normal shape for an
\emph{Education} or \emph{Certificates} section, where the heading carries the
degree, institution, and dates --- renders its heading and contributes nothing
further, and the boundary after it stays a legal page break.

The body is read as an argument, so catcode-sensitive content cannot appear
directly inside an entry. \cmd{verb}, \textsf{listings} environments, and
anything else that depends on rescanning characters must be wrapped in a macro
defined outside the entry, or placed outside it.

\subsection{\env{CDossierItemize}}
\label{sec:itemize}

\begin{latexcode}
\begin{CDossierItemize}
  \item First accomplishment.
  \item Second accomplishment.
\end{CDossierItemize}
\end{latexcode}

A résumé-appropriate itemized list with controlled indentation and spacing.
Prefer it over globally redefining \env{itemize}.

The space above the list is \cmd{CDossierRecordListEdgeAboveSkip} and the space
below is \cmd{CDossierRecordListEdgeBelowSkip}. Each token is the complete gap
at its end of the list, and both collapse with the adjacent block's own spacing
rather than adding to it. The CV's publication list uses the same pair.

\cmd{CDossierRecordListEdgeAboveSkip} additionally carries a \textbf{lower bound
of \texttt{0.25}}, which is an extraction constraint rather than a design
preference. The entry heading sets its dates and location in a right-hand
column, and a geometric extractor keeps that column with its entry only while
the entry's vertical band stays distinct from the list beneath it. Below
\texttt{0.25} the two merge and the dates extract after the bullets --- on the
résumé and the CV alike, at every supported body size, and on tagged and
untagged output equally. Overriding this token below \texttt{0.25} therefore
breaks reading order; \opt{entrymeta=inline} is the supported way to go lower.

When the list crosses a page boundary it is never split so that a single item
stands alone on either side; a list longer than a page still breaks normally.

\subsection{Manual publications}
\label{sec:publications}

Manual publications require no bibliography package and no Biber run:

\begin{latexcode}
\CDossierSection{Publications}

\begin{CDossierPublications}
  \CDossierPublication{
    authors = {Ada Lovelace and Jane Example},
    title   = {A Demonstration Article},
    venue   = {Journal of Examples},
    date    = {2026},
    doi     = {10.9999/example.2026.1}
  }
\end{CDossierPublications}
\end{latexcode}

\env{CDossierPublications} creates a numbered list in source order and uses the
shared list spacing and label separation. \cmd{CDossierPublication} is valid
only inside that environment, and takes these keys:

\begin{longtable}{@{}llL{0.50\textwidth}@{}}
\toprule
\textbf{Key} & \textbf{Required} & \textbf{Purpose} \\
\midrule
\endhead
\key{authors} & Yes & Display-order author list \\
\key{title} & Yes & Publication title \\
\key{venue} & No & Journal, book, conference, or publisher \\
\key{date} & No & Year or display date \\
\key{doi} & No & DOI value or complete DOI URL \\
\key{url} & No & Fallback public URL \\
\key{note} & No & Short status or contribution note \\
\bottomrule
\end{longtable}

Missing optional values collapse cleanly. When both \key{doi} and \key{url} are
present, the DOI is the displayed link and the URL is the fallback. Each entry
renders in this order: authors, italic title, the comma-joined present
\key{venue} and \key{date} values, \key{note}, and the preferred visible
identifier. Sentence punctuation is emitted only around present groups, so
absent optional fields leave no stray separators.

\subsubsection{Numbering across grouped lists}

The environment takes one optional key:

\begin{longtable}{@{}lllL{0.40\textwidth}@{}}
\toprule
\textbf{Key} & \textbf{Values} & \textbf{Default} & \textbf{Purpose} \\
\midrule
\endhead
\key{numbering} & \opt{restart}, \opt{continue} & \opt{restart} & Whether this list numbers from 1 or carries on from the previous list \\
\bottomrule
\end{longtable}

Each list is self-contained by default, so a section carrying several groups
numbers each from 1. That reads correctly when a number is cited together with
its group --- ``journal article 1''. A document whose numbers are cited on their
own, from a cover letter or a grant form, wants one sequence across the whole
document instead, and asks for it per list:

\begin{latexcode}
\CDossierSubsection{Conference Papers}
\begin{CDossierPublications}[numbering = continue]
  ...
\end{CDossierPublications}
\end{latexcode}

\opt{continue} resumes from the final number of the \emph{preceding} list,
wherever that list was: a subsection heading and a section rule are equally
invisible to it, since the sequence belongs to the document rather than to the
heading above it. A \opt{restart} list in between therefore resets what a later
\opt{continue} picks up. Setting \key{numbering} on the first list of a document
has no effect under either value.

\subsection{\cmd{CDossierLink}}
\label{sec:link}

\begin{latexcode}
\CDossierLink{example.org/ada/portfolio}
\CDossierLink{https://example.org/programs/2024/final-report/notes.html}
\end{latexcode}

The supported way to put a link in body text --- a bullet, an entry, a letter
paragraph, a statement paragraph. The argument is both the text that appears on
the page and, once normalized, the link target: a value carrying no scheme gains
\texttt{https://}, exactly as \key{website} does, and a value that has one keeps
it. The displayed text is the argument as written, never the normalized target.

\cmd{CDossierLink} takes no bare identifier: body text has no service to expand
one against (\S\ref{sec:webkeys}).

Use it for every address in running text. Typing the URL as plain text does not
work, and neither does \cmd{url}:

\begin{longtable}{@{}lL{0.66\textwidth}@{}}
\toprule
\textbf{Form} & \textbf{What a long address does} \\
\midrule
\endhead
plain text & no URL break points, so \TeX{} either overruns the margin or hyphenates the address --- and the inserted hyphen travels into the pasted URL \\
\cmd{url} & breaks after punctuation only, so a long segment carrying none stays one unbreakable token \\
\cmd{CDossierLink} & may break after any character, with no inserted hyphen \\
\bottomrule
\end{longtable}

\cmd{CDossierEmergencyStretch} cannot rescue either of the first two: it
redistributes a line's interword glue, and a line holding one over-wide token
has none to redistribute.

\cmd{CDossierLink} extends the URL break set with the ASCII letters and digits
for the duration of one link, so the address may break anywhere rather than only
after punctuation, and a break point inserts a penalty and no discretionary
hyphen --- so a wrapped address still copies and pastes back character for
character. The extension is local: the contact line and the bibliography keep
the punctuation-only breaks they were calibrated with.

The argument is read as ordinary text (\S\ref{sec:escaping}). A query string is
the common case:

\begin{latexcode}
\CDossierLink{https://example.org/search?q=studies\&profile=advanced}
\end{latexcode}

Keep the address on one source line. A URL cannot contain a literal space, so a
space in the argument --- in practice always a line break inside it --- is
removed from the link target, with a warning naming the repaired target. Without
that repair the page would look correct while the link went somewhere else.
Write \texttt{\%20} if the space is genuinely part of the address.

A \cmd{CDossierLink} is never underlined under either \opt{medium}: its visible
text is the address, so it already announces itself.

\subsection{\cmd{CDossierPlainLinks}}

\begin{latexcode}
{\CDossierPlainLinks
 Read the \href{https://example.org/report}{annual report}.}
\end{latexcode}

A declaration that turns off link decoration for the rest of the enclosing
group. Under \opt{medium=print} it changes nothing, because nothing is decorated
there; under \opt{medium=screen} an \cmd{href} inside its scope is set without
the underline while still emitting its link annotation.

Use it where a link's visible text is already the address, or where a passage
carries so many links that ruling each one would be noise. The toolkit applies
it to its own address-as-text links --- the contact line, \cmd{CDossierLink},
the ORCID iD, the manual publication identifiers, and the whole bibliography ---
so those need nothing from a document.

There is no matching command to turn decoration back on. Scope it with a group,
as above, or with the environment it belongs to.

\subsection{Field accessors}

These are intended mainly for advanced users and shared components.

\cmd{CDossierPrintField}\texttt{\{name\}} prints the stored value of a profile
field. For an absent optional field it prints nothing; for an unknown field name
it produces an actionable error.

\cmd{CDossierIfFieldTF}\texttt{\{phone\}\{present\}\{absent\}} executes the
first branch when the field exists and is nonblank, and the second otherwise.
Both branches are long, so they may contain full paragraphs.

\cmd{CDossierFieldValue}\texttt{\{email\}} expands to the stored value, or to
nothing when the field is absent. It does not typeset and does not validate the
field name.

Expandability is what makes the last of these the correct choice anywhere the
value is written into the PDF rather than onto the page --- a heading's recorded
title, or any other PDF string. Such values are purified first, and purification
expands what it can and discards what it cannot: \cmd{CDossierPrintField} is
protected, so it does not survive as a value but its argument text does, which
records a plausible wrong string rather than failing. Nothing on the rendered
page reveals that.

\subsection{Hyperlink behaviour}

When present, \key{email} uses a \texttt{mailto:} link, and \key{website},
\key{linkedin}, \key{github}, and \key{scholar} use web links. Printed text
stays readable in monochrome, and URLs are breakable rather than extending
beyond the margin. A displayed URL is not assumed to include a protocol.

\paragraph{Without \textsf{hyperref}.} Every link the toolkit emits is emitted
only when \textsf{hyperref} has supplied \cmd{href}. Without it each one
degrades to its visible address as plain text, so the identifier stays readable,
selectable, and extractable; nothing in the PDF is clickable.

All four classes load \textsf{hyperref}, so a document built on one of them is
never in this state. It is reachable by loading
\file{careerdossier-components} into another class, and there the degradation
is otherwise invisible --- the page looks finished. The package therefore warns
once per document, at the first link that degrades, naming the cause and the
effect. Once, not once per field: a contact line with five linked fields has a
single cause.

\section{Optional BibLaTeX and Biber}
\label{sec:biblatex}

The CV class does not load \textsf{biblatex}. Opt in explicitly:

\begin{latexcode}
\documentclass{careerdossier-cv}
\usepackage{careerdossier-biblatex}

\addbibresource{publications.bib}
\CDossierHighlightAuthor{
  family = {Lovelace},
  given  = {Ada}
}

\begin{document}
\nocite{*}
\printbibliography[title={Publications}]
\end{document}
\end{latexcode}

The integration package loads and configures \textsf{biblatex}; standard
BibLaTeX commands remain the public resource-selection and printing interface.
Its one supported profile is fixed:

\begin{plaincode}
backend=biber
style=numeric
sorting=ydnt
\end{plaincode}

Entries are numbered, sorted year-descending / name / title, and show at most
one preferred public identifier, in the order DOI, e-print, URL. The package
uses monochrome link styling and does not redefine unrelated document lists or
headings globally.

The gap between bibliography entries follows the same calibrated list token as
every other list in the CV, so a \textsf{biblatex} list and a manual publication
list share one rhythm at every \opt{fontsize}. The horizontal gap between an
entry number and its entry follows the shared list label token for the same
reason. BibLaTeX's own default places that gap at roughly one em, which is the
width at which a plain-text extractor stops treating the numbers as labels and
starts treating them as a separate left-hand column, emitting them as a block
ahead of the entry text; the shared token is half the body size and clears that
threshold at every supported \opt{fontsize}.

URLs printed in a bibliography carry no stretch at their break points, so a
justified line ending in a URL spreads its word spaces rather than the URL
itself --- a stretched URL extracts as separate tokens and is then neither
searchable nor copyable. Rigid glue rules that out at any measure or URL length,
and the profile therefore also enables BibLaTeX's three URL break penalties,
which are disabled by default: without them a long URL would overrun the margin
instead. Two visible results follow, both intended. A URL breaks where the line
ends rather than at the nearest earlier punctuation, so an address may wrap in
the middle of a run of letters or digits; and a line holding nothing but URL is
set ragged-right and reported underfull, since it has no glue to justify with. A
wrapped address still concatenates back exactly.

\cmd{CDossierHighlightAuthor} may be repeated for spelling or initial variants.
It bolds an exact BibLaTeX-parsed family/given-name pair in the bibliography and
does not alter citations. Both keys are required; an incomplete declaration
produces an actionable package error.

Loading \file{careerdossier-biblatex} when \textsf{biblatex} is unavailable
reports the missing optional dependency and explains that you may either install
BibLaTeX and Biber or use \env{CDossierPublications}. A CV that does not load
the integration package builds without either.

\section{PDF document metadata}
\label{sec:metadata}

Each class derives the PDF's document metadata from the shared profile. Nothing
needs to be called; the derived values are applied automatically at
\cmd{begin}\texttt{\{document\}}.

\begin{longtable}{@{}L{0.16\textwidth}L{0.18\textwidth}L{0.58\textwidth}@{}}
\toprule
\textbf{PDF field} & \textbf{Derived from} & \textbf{Value} \\
\midrule
\endhead
\texttt{/Title} & \key{name} & \texttt{Résumé – <name>}, \texttt{Cover Letter – <name>}, \texttt{Curriculum Vitae – <name>}, \texttt{Academic Cover Letter – <name>}, or the statement's title \\
\texttt{/Author} & \key{name} & \texttt{<name>} \\
\texttt{/Lang} & fixed & \texttt{en} \\
\bottomrule
\end{longtable}

The document type is part of the title so that a résumé and a cover letter built
from one profile are distinguishable in a viewer's tab bar, in document
properties, and in a file manager. The separator is an en dash, and the derived
values are the same whether or not the document opts into tagging.

When \key{name} is absent, \texttt{/Title} and \texttt{/Author} are left unset;
the header commands already error on a missing \key{name}, and metadata does not
add a second diagnostic.

\subsection{\TeX{} input ligatures in a derived value}

Text you supply is \TeX{} input, so the usual input ligatures apply to it on the
way into the PDF. A \key{name} of \texttt{\{Ada Lovelace-{}-Byron\}} reaches
\texttt{/Title} and \texttt{/Author} with an en dash, and the same holds for a
statement's \key{title}:

\begin{longtable}{@{}L{0.30\textwidth}L{0.60\textwidth}@{}}
\toprule
\textbf{You type} & \textbf{The PDF carries} \\
\midrule
\endhead
\texttt{-{}-} & – (en dash, U+2013) \\
\texttt{-{}-{}-} & — (em dash, U+2014) \\
\texttt{!`} & ¡ \\
\texttt{?`} & ¿ \\
\bottomrule
\end{longtable}

Nothing else is converted. In particular the two ASCII quote sequences are
\textbf{not} turned into curly quotes in a metadata string --- they arrive as
the characters you typed. Type the character itself if you want one:
\texttt{\{Ada Lovelace–Byron\}} in a UTF-8 source works and is unambiguous.

This holds on both build paths. Only values this package derives are covered;
see below for a value you set yourself.

\subsection{\texttt{/DisplayDocTitle}}

A \texttt{/Title} in the file is not what a viewer puts in its window title, tab
bar, or recent-documents list. A viewer uses the filename unless the PDF's
viewer preferences ask otherwise, so the classes also request
\texttt{/DisplayDocTitle true}. Without it the derived title is present and
unused --- the case PDF/UA-2 clause 8.11.2 and WCAG 2.1 AA 2.4.2 (Page Titled)
are about. It is requested on both build paths and changes nothing about
rendering, extracted text, or the structure tree.

\subsection{Overriding the derived metadata}

Set any field yourself with \cmd{hypersetup} and it is used unchanged:

\begin{latexcode}
\documentclass{careerdossier-resume}

\hypersetup{
  pdftitle  = {Ada Lovelace — Data Scientist, Analytical Engine Division},
  pdfauthor = {A. Lovelace},
  pdflang   = {en-GB}
}

\CDossierSetup{ name = {Ada Lovelace} }
\end{latexcode}

A field you set is never overwritten, and the order does not matter --- the
\cmd{hypersetup} may appear before or after \cmd{CDossierSetup}. Fields you do
not set are still derived, so overriding \texttt{pdftitle} alone leaves
\texttt{/Author} and \texttt{/Lang} in place. This holds whether or not you opt
into tagging.

The language has a second route, and it wins too:

\begin{latexcode}
\DocumentMetadata{ lang = de, tagging = on }
\documentclass{careerdossier-resume}
\end{latexcode}

\cmd{DocumentMetadata} records the language for the \LaTeX{} kernel, which
writes \texttt{/Lang} itself, and the derived \texttt{en} stands aside. It always
settles on a language --- it defaults to \texttt{en} when you give no
\key{lang} key --- so on that path the kernel owns \texttt{/Lang} outright.

Declaring a language does not make the document that language.
CareerDossierTeX is English-only; \key{lang} sets the declaration a screen
reader reads, not hyphenation, fonts, or any fixed string.

\texttt{pdfdisplaydoctitle} is the one setting that works by ordering rather
than by detection, because \textsf{hyperref} gives a boolean no state that
distinguishes ``not set'' from ``set to false''. The classes request it while
\textsf{hyperref} is still loading, so any value in your preamble is later and
wins:

\begin{latexcode}
\hypersetup{ pdfdisplaydoctitle = false }   % show the filename instead
\end{latexcode}

The one place that does not reach is a value passed as a package option, which
\textsf{hyperref} processes before the classes can be asked anything. Set it
with \cmd{hypersetup} in the preamble instead.

Other \textsf{hyperref} metadata (\texttt{pdfsubject}, \texttt{pdfkeywords},
\dots) is untouched.

\paragraph{One divergence to know about.} A value you set yourself does not get
the input-ligature handling described above, because this package never rewrites
it. It follows \textsf{hyperref}'s own behaviour, and that behaviour differs
between the two build paths: \textsf{hyperref} converts on the default path,
while under \cmd{DocumentMetadata} your value goes to the \LaTeX{} kernel
instead and arrives as typed. So

\begin{latexcode}
\hypersetup{ pdftitle = {Résumé -- Ada Lovelace} }
\end{latexcode}

gives an en dash without \cmd{DocumentMetadata} and two hyphens with it. The
same applies to \texttt{pdfsubject} and \texttt{pdfkeywords}, and to a plain
\textsf{article} with \textsf{hyperref} and no CareerDossierTeX file loaded ---
which is what places the cause upstream rather than in these classes. Type the
character you want, and both paths agree on every field at once.

\section{Page furniture}
\label{sec:furniture}

All four classes use the same page-furniture policy. Under the default
\opt{medium=print}:

\begin{itemize}
  \item a one-page document has no running header and no folio;
  \item page one of a multi-page document has only a centred
        \texttt{Page N of M} folio;
  \item pages two and later add a centred \texttt{<name> -- <document label>}
        running header above the body and retain the folio;
  \item the résumé label is \texttt{Résumé}, the CV label is
        \texttt{Curriculum Vitae}, both letter families use
        \texttt{Cover Letter}, and a statement uses its independently
        configurable short running title.
\end{itemize}

Under \opt{medium=screen} neither is emitted, on any page and at any page count.

Furniture uses the sans-serif furniture step of the calibrated type scale and
the monochrome text token. It stays print-safe and is marked as a layout
artifact when tagged output is active. Which furniture is emitted is the only
user-configurable part of the policy; its typography, wording, and placement are
not.

The header and the folio sit in the vertical centre of the top and bottom
margins respectively, at every \opt{fontsize} $\times$ \opt{margin} combination
and on both paper sizes. Their distance from the paper edge therefore follows
the selected margin preset. This affects only where the furniture sits inside
the margin; the text block is unchanged.

Single-page suppression and the \texttt{of M} total use the last absolute page
recorded in the auxiliary file, which is why a first build from a clean tree
shows a provisional folio.

\section{Pagination}
\label{sec:pagination}

The résumé and CV classes state where a page may break rather than leaving every
break to \LaTeX's defaults:

\begin{itemize}
  \item a section heading is placed only when it has room to introduce
        something;
  \item an entry heading stays with its own second line and with the first line
        of its body;
  \item a bullet list is never split leaving one item alone on either side.
\end{itemize}

The policy uses page-break penalties, not boxing, so material that genuinely
does not fit still breaks: an entry or bullet list longer than a page paginates
normally rather than overflowing.

\subsection{\cmd{CDossierSectionNeedLines}}

The section heading's keep, expressed as a bound. A section heading is placed
only when at least this many lines of body leading remain on the page;
otherwise the page breaks before the heading and the heading opens the next one.
It defaults to \texttt{4}, which is what the heading needs to introduce
anything: the heading and its rule, the two lines of the first entry heading,
and that entry's first line of body.

It is a count of lines rather than a dimension, because it is a statement about
content --- the same four lines at every \opt{fontsize}.

\begin{latexcode}
\int_set:Nn \CDossierSectionNeedLines { 6 }
\end{latexcode}

\subsection{\cmd{CDossierSubsectionNeedLines}}

The same bound for \cmd{CDossierSubsection}, defaulting to \texttt{3}. It is
smaller because a subsection commits less to the page: its own heading line and
the two lines of the first entry heading, with no rule between them. It does not
additionally require that entry's first line of body, because the entry heading
already carries its own keep down to it.

A subsection that opens its section directly is exempt from this bound: the
section's own requirement already placed that material, and testing it a second
time could only break the page between a section heading and the subsection
immediately under it --- the stranded heading the bound exists to prevent.

\subsection{Typographic page-break penalties}
\label{sec:penalties}

All four classes set \cmd{brokenpenalty}, \cmd{clubpenalty}, and
\cmd{widowpenalty} from three named tokens ---
\cmd{CDossierBrokenPenalty}, \cmd{CDossierClubPenalty}, and
\cmd{CDossierWidowPenalty} --- through one command,
\cmd{CDossierApplyPageBreakPenalties}. A hyphenated word is therefore never
split across a page break, and no single line of a paragraph is stranded alone
at the foot of one page or the head of the next.

All three default to \texttt{10000} --- forbidding the break --- across every
family. This is safe because all four classes are \cmd{raggedbottom}:
forbidding a club or widow break only closes off the first and last line of a
paragraph as a break point, and every interior line break remains available, so
a page can still break inside an over-long paragraph rather than overflow.

These tokens are \textbf{booleans, not dials}. Under \cmd{raggedbottom} a page
being built carries no stretch, so every break candidate short of the goal
scores the worst possible badness; all candidates tie, the page builder keeps the
last one that fits, and the penalty's magnitude never enters the comparison.
Only a value of \texttt{10000} or more changes anything, by removing the
breakpoint from consideration altogether. Measured across every letter and
statement fixture at \texttt{10000, 4000, 1000, 300, 150, 50, 0}, the results
form a step function with a single step: \texttt{9999} paginates identically to
\texttt{0}. Setting either token to a ``discounted'' value does not soften the
rule, it switches it off.

The cost of that rule was measured and accepted: each fixture that gains page
fill by permitting these breaks gains a club or widow line with it, one for one.
There is no value, and no per-family split, that buys the space without the
defect, so the toolkit keeps the space.

The letter and statement classes otherwise remain continuous prose with no
further structural page-break policy.

\subsection{Hyphenation}

\cmd{hyphenpenalty} and \cmd{exhyphenpenalty} keep \TeX's defaults of
\texttt{50} in all four families. No token sets them, and that is a decision
rather than an omission: raising the penalty was measured and rejected. Across
the shipped examples, \texttt{500} removes 18 hyphens and creates 10 lines
looser than badness 99 --- roughly a one-for-one trade, paid to fix a problem
the record classes do not have. Forbidding hyphenation entirely produces
overfull boxes.

A document that wants different behaviour can set either parameter in its
preamble; both are ordinary \TeX{} parameters and nothing in the toolkit
overrides them.

\section{Design tokens}
\label{sec:tokens}

\subsection{Vertical spacing}

Every vertical boundary in the four classes is owned by exactly one token, and
each token is the complete gap a reader measures. The names follow
\cmd{CDossier<Family><Scope><Position>Skip}.

\begin{longtable}{@{}lL{0.70\textwidth}@{}}
\toprule
\textbf{Family} & \textbf{Tokens} \\
\midrule
\endhead
Shared (all four) & \cmd{CDossierSharedHeaderNameGapSkip},
  \cmd{CDossierSharedHeaderMetaGapSkip} \\
Record (résumé, CV) & \cmd{CDossierRecordHeaderBelowSkip},
  \cmd{CDossierRecordSectionAboveSkip},
  \cmd{CDossierRecordSectionRuleGapSkip},
  \cmd{CDossierRecordSectionBelowSkip},
  \cmd{CDossierRecordSubsectionAboveSkip},
  \cmd{CDossierRecordSubsectionBelowSkip},
  \cmd{CDossierRecordEntryAboveSkip},
  \cmd{CDossierRecordEntryGapSkip},
  \cmd{CDossierRecordListEdgeAboveSkip},
  \cmd{CDossierRecordListEdgeBelowSkip},
  \cmd{CDossierRecordItemSepSkip},
  \cmd{CDossierRecordParSkip} \\
Prose (statement) & \cmd{CDossierProseHeaderBelowSkip},
  \cmd{CDossierProseSectionAboveSkip},
  \cmd{CDossierProseSectionBelowSkip},
  \cmd{CDossierProseSubsectionAboveSkip},
  \cmd{CDossierProseSubsectionBelowSkip},
  \cmd{CDossierProseParSkip} \\
Letter & \cmd{CDossierLetterHeaderBelowSkip},
  \cmd{CDossierLetterParSkip},
  \cmd{CDossierLetterRecipientLineGapSkip},
  \cmd{CDossierLetterBlockSkip},
  \cmd{CDossierLetterBodyAboveSkip},
  \cmd{CDossierLetterBodyBelowSkip},
  \cmd{CDossierLetterSignatureGapSkip} \\
\bottomrule
\end{longtable}

Only the two gaps \emph{inside} the header block are shared; the gap
\emph{below} it is owned by one token per family. All three ship at the same
ratio, so setting all three reproduces a single shared token, and setting one
moves only that family.

Three rules govern how a token reaches the page, and overriding a token without
them in mind is the usual reason an override appears to do nothing.

\begin{itemize}
  \item \textbf{Boundaries take the maximum of the adjacent claims, never their
        sum.} Raising a token past its neighbour is what makes it visible;
        lowering it below its neighbour makes it inert, and neither the output
        nor the log will say so. \cmd{CDossierRecordEntryGapSkip} in particular
        is a \emph{floor} for the entry-heading-to-body boundary, which a bullet
        list overrides with \cmd{CDossierRecordListEdgeAboveSkip}.
  \item \textbf{One token carries a constraint from outside the type scale.}
        \cmd{CDossierRecordListEdgeAboveSkip} may not go below \texttt{0.25}
        without breaking the extraction order of the entry heading's dates
        column (\S\ref{sec:itemize}).
  \item \textbf{A paragraph boundary also contributes \cmd{parskip}, and the
        class has already accounted for it.} Where a token sits at such a
        boundary the class emits it minus \cmd{parskip}, so the value you set is
        still the gap a reader measures --- do not subtract the paragraph gap
        yourself.
\end{itemize}

The calibrated \emph{values} are not stable API before \texttt{v0.10.0}; the
token names and the boundaries they own are.

\subsection{Line breaking}

\cmd{CDossierEmergencyStretch} is the extra flexibility \TeX{} may distribute
among a paragraph's interword glue when it can find no set of breakpoints within
\cmd{tolerance}. It is a fixed multiple of the body size, and it derives from
the body size rather than the body leading because the quantity is horizontal.
The ratio, and the points it resolves to at each \opt{fontsize}, are in
\file{docs/ARCHITECTURE.md} with the rest of the calibrated values. It does
not scale with the measure, and deliberately so: deriving it from the line width
instead was measured and rejected, because the pool a paragraph needs tracks
neither quantity.

All four classes set \cmd{emergencystretch} from it, unconditionally and
identically. Raising it lets a paragraph with few break points set looser
instead of overfull; lowering it to \texttt{0pt} restores \TeX's default
behaviour, which is to allow the overfull box. Neither direction affects a
paragraph that already sets within \cmd{tolerance}, because
\cmd{emergencystretch} is consulted only in a third line-breaking pass that such
a paragraph never reaches.

\subsection{Typography roles}

Each role names one place in a document, and describes meaning rather than a
particular font family, weight, or size:

\begin{longtable}{@{}lL{0.60\textwidth}@{}}
\toprule
\textbf{Role} & \textbf{Where it applies} \\
\midrule
\endhead
\cmd{CDossierNameStyle} & the profile name at the head of every document \\
\cmd{CDossierHeadlineStyle} & the optional headline below the name \\
\cmd{CDossierSectionStyle} & every section heading, in every class; also, at a
  smaller step of the size scale, the statement's subsection headings, and at a
  larger one its document title \\
\cmd{CDossierEntryTitleStyle} & the title line a \env{CDossierEntry} renders, in
  the résumé and the CV \\
\cmd{CDossierSubjectStyle} & a cover letter's subject line \\
\cmd{CDossierBodyStyle} & running prose, and every line set as prose: the
  contact line, an entry's organization, and the letter's date, recipient block,
  salutation, closing, and signature name \\
\cmd{CDossierMutedStyle} & the de-emphasised runs \opt{muted} governs \\
\cmd{CDossierLinkStyle} & author-written anchor text under \opt{medium=screen};
  the identity under \opt{medium=print} and inside \cmd{CDossierPlainLinks} \\
\bottomrule
\end{longtable}

\cmd{CDossierEntryTitleStyle} and \cmd{CDossierSubjectStyle} resolve to the same
shape today. That is two independent decisions agreeing, not one role with two
names: redefining either leaves the other alone. Do not define one in terms of
the other.

\cmd{CDossierMutedStyle} and \cmd{CDossierLinkStyle} are published by the
components package rather than the typography package, because each resolves to
something the typography package does not own --- a colour and a rule
respectively. Both are available in every document class. A document that wants
a different link affordance redefines \cmd{CDossierLinkStyle}.

Their visual definitions may evolve before \texttt{v0.10.0}.

\subsection{Colour}

The monochrome theme exposes four semantic tokens:

\begin{latexcode}
\CDossierTextColor
\CDossierMutedColor
\CDossierRuleColor
\CDossierLinkColor
\end{latexcode}

\cmd{CDossierMutedColor} is gray 0.30, measured at 8.52:1 against white under
the WCAG 2.1 relative-luminance formula. It is what \opt{muted=gray} and
\opt{muted=both} render de-emphasised runs in; under \opt{muted=plain} or
\opt{muted=italic} nothing uses it.

Do not rely on the underlying colour names or values as stable API before
\texttt{v0.10.0}.

\section{Tagged structure}
\label{sec:tagging}

Tagged output is opt-in and off by default. It is enabled with the \LaTeX{}
kernel's \cmd{DocumentMetadata}, which must appear \textbf{before}
\cmd{documentclass}:

\begin{latexcode}
\DocumentMetadata{lang=en, tagging=on}
\documentclass{careerdossier-resume}
\end{latexcode}

This introduces no CareerDossierTeX class option and no public command. The
tagging interface is the kernel's, not this package's.

When tagging is active the classes expose section headings, lists, paragraphs,
and links as structure, and mark decorative rules, contact separators, and
running page furniture as layout artifacts. When it is not active, output is
unchanged from the untagged path. Under tagging, \cmd{CDossierSubsection} is a
depth-3 heading role-mapped to \texttt{/H3}, opening a depth-3 division nested
inside its section's.

\paragraph{Scope of the claim.} Tagged output is a tested preview for five
fixture profiles --- industry résumé, industry letter, academic CV, academic
letter, and statement. It is \textbf{not} a PDF/UA, WCAG, ATS, or general
accessibility conformance claim. The macOS VoiceOver reading-order pass covers
four of the five; the statement fixture and Windows/NVDA remain
screen-reader-unverified.

\section{Errors and warnings}

\subsection{Errors}

Compilation stops for:

\begin{itemize}
  \item use under an unsupported engine;
  \item an unsupported class-option value;
  \item a missing required \key{name} when rendering identity content;
  \item an unknown public field or label name;
  \item a manual publication used outside \env{CDossierPublications};
  \item a manual publication missing \key{authors} or \key{title};
  \item an unknown manual-publication, publication-list, or preferred-author
        key;
  \item an unsupported \key{numbering} value;
  \item a preferred-author declaration missing \key{family} or \key{given};
  \item loading \file{careerdossier-biblatex} when BibLaTeX is unavailable;
  \item malformed key-value input that cannot be interpreted safely.
\end{itemize}

A missing-BibLaTeX diagnostic names BibLaTeX and Biber, recommends building with
latexmk after installation, and points to \env{CDossierPublications} as the
dependency-free alternative. A missing Biber executable is reported by the
standard toolchain rather than by a separate preflight.

\subsection{Warnings}

Warnings are raised for an unusually long contact field, a URL that cannot be
normalized, an \key{orcid} value that does not have the shape of an ORCID iD,
metadata overwritten by a later setup call, fields accepted but not displayed by
the active class, and links degraded to plain text because \textsf{hyperref} is
absent. Each explains the likely effect and the corrective action.

\section{Complete examples}
\label{sec:examples}

\subsection{Résumé}

\begin{latexcode}
\documentclass{careerdossier-resume}

\CDossierSetup{
  name     = {Ada Lovelace},
  headline = {Data Scientist},
  email    = {ada@example.com},
  location = {London, UK}
}

\begin{document}

\MakeCDossierHeader

\CDossierSection{Experience}

\begin{CDossierEntry}[
  title        = {Data Scientist},
  organization = {Example Organization},
  location     = {Toronto, Ontario},
  dates        = {2024--Present}
]
  \begin{CDossierItemize}
    \item Developed a reproducible analytical workflow.
  \end{CDossierItemize}
\end{CDossierEntry}

\end{document}
\end{latexcode}

\subsection{Cover letter}

\begin{latexcode}
\documentclass{careerdossier-letter}

\CDossierSetup{
  name     = {Ada Lovelace},
  headline = {Data Scientist},
  email    = {ada@example.com},
  location = {London, UK}
}

\CDossierLetterSetup{
  recipient-name         = {Hiring Manager},
  recipient-organization = {Example Organization},
  subject                = {Application for the Data Scientist Position}
}

\begin{document}

\MakeCDossierLetterhead

I am writing to apply for the Data Scientist position.

\MakeCDossierClosing

\end{document}
\end{latexcode}

An academic cover letter is the same document with
\cmd{documentclass}\texttt{[family=academic]}.

\subsection{Academic CV}

\begin{latexcode}
\documentclass{careerdossier-cv}

\CDossierSetup{
  name    = {Ada Lovelace},
  email   = {ada@example.com},
  scholar = {kukA0LcAAAAJ},
  orcid   = {0000-0002-1825-0097}
}

\begin{document}

\MakeCDossierHeader

\CDossierSection{Academic Appointments}

\begin{CDossierEntry}[
  title        = {Assistant Professor},
  organization = {Example University},
  location     = {Toronto, Ontario},
  dates        = {2024--Present}
]
  \begin{CDossierItemize}
    \item Research and teaching summary.
  \end{CDossierItemize}
\end{CDossierEntry}

\CDossierSection{Publications}

\begin{CDossierPublications}
  \CDossierPublication{
    authors = {Ada Lovelace and Jane Example},
    title   = {A Demonstration Article},
    venue   = {Journal of Examples},
    date    = {2026},
    doi     = {10.9999/example.2026.1}
  }
\end{CDossierPublications}

\end{document}
\end{latexcode}

Education, appointments, research, teaching, grants, awards, presentations, and
service all use \env{CDossierEntry}; there is no separate command per section
type.

\subsection{Statements}

The default interest type needs only \key{name} and \key{email}:

\begin{latexcode}
\documentclass{careerdossier-statement}
\CDossierSetup{name={Ada Lovelace}, email={ada@example.com}}
\begin{document}
\MakeCDossierStatementHeader
This statement introduces work and interests without a type-specific contract.
\end{document}
\end{latexcode}

Research requires \key{affiliation}:

\begin{latexcode}
\documentclass[type=research]{careerdossier-statement}
\CDossierSetup{
  name={Ada Lovelace}, email={ada@example.com},
  affiliation={Example University}, orcid={0000-0002-1825-0097}
}
\begin{document}
\MakeCDossierStatementHeader
My research develops reliable methods for computational inquiry.
\end{document}
\end{latexcode}

Artist requires \key{website}:

\begin{latexcode}
\documentclass[type=artist]{careerdossier-statement}
\CDossierSetup{
  name={Ada Lovelace}, email={ada@example.com},
  website={portfolio.example.com}
}
\begin{document}
\MakeCDossierStatementHeader
My practice examines the relationship between material and computation.
\end{document}
\end{latexcode}

Purpose may carry both an application context and an identifier:

\begin{latexcode}
\documentclass[type=purpose]{careerdossier-statement}
\CDossierSetup{name={Ada Lovelace}, email={ada@example.com}}
\CDossierStatementSetup{
  application-context = {Application to the MSc in Computational Science},
  application-id      = {12345678}
}
\begin{document}
\MakeCDossierStatementHeader
I seek advanced study in reliable scientific computing.
\end{document}
\end{latexcode}

The remaining types --- \opt{teaching}, \opt{teaching-philosophy}, and
\opt{diversity} --- follow the interest example, differing only in the title
their type selects.

\section{Further reading}

These live in the repository rather than in this manual, because each is about
the project rather than about writing a document with it.

\begin{longtable}{@{}lL{0.60\textwidth}@{}}
\toprule
\textbf{Document} & \textbf{Subject} \\
\midrule
\endhead
\file{README.md} & overview, installation, and a short quick start \\
\file{CHANGELOG.md} & release history and user-visible changes \\
\file{docs/MIGRATION.md} & what changed between releases, and what needs a source edit \\
\file{docs/ARCHITECTURE.md} & module boundaries, internal design, and the calibrated ratios behind the tokens \\
\file{docs/ATS-EXTRACTION.md} & design guidance for ATS-safe output and text extraction \\
\file{docs/ROADMAP.md} & release phases, non-goals, and what is deferred \\
\file{docs/TESTING.md} & the test suites and coverage expectations \\
\file{CONTRIBUTING.md} & issue, branch, commit, and pull-request workflow \\
\bottomrule
\end{longtable}

Report a defect or propose a feature at
\url{https://github.com/amirhs1/CareerDossierTeX/issues}.

\end{document}
